% =====================================================================
%  Two machine-checked refutations in the Fresco metatheory.
%
%  BUILD (from mechanization/docs/):
%      TEXINPUTS=../../paper: pdflatex refutations.tex
%
%  Standalone document: it has its own preamble and body, and reuses the
%  paper's notation by \input-ing paper/macro.tex, so every judgement
%  below is typeset by the paper's own macros.
% =====================================================================
\documentclass[11pt,a4paper]{article}

\usepackage[utf8]{inputenc}
\usepackage[T1]{fontenc}
\usepackage{amsmath,amssymb,mathtools,amsthm}
\usepackage{graphicx}
\usepackage{thmtools}
\usepackage{mathpartir}
\usepackage{stmaryrd}
\usepackage{xspace}
\usepackage{url}
\usepackage{xcolor}
\usepackage{array}
\usepackage{tikz}
\usepackage[normalem]{ulem}
\usepackage{enumerate}
\usepackage{enumitem}
\usepackage{booktabs}
\usepackage{halloweenmath}
\usepackage{trimclip}
\usepackage{scalerel}[2016/12/29]
\usepackage[old]{old-arrows}
\usepackage[ruled, noend]{algorithm2e}
\usepackage{listings}
\usepackage{pifont}
\usepackage{fontawesome}
\usepackage[margin=1in]{geometry}
\usepackage{hyperref}
\usepackage{cleveref}

\input{macro}

\newtheorem*{thmpaper}{Theorem 4.6}
\newtheorem*{lempaper}{Lemma E.17}
\newtheorem*{defpaper}{Definition (Problems and solutions of subtyping)}
\theoremstyle{definition}
\newtheorem{claimc}{Claim}
\crefname{claimc}{Claim}{Claims}
\Crefname{claimc}{Claim}{Claims}

\newcommand{\coq}[1]{\texttt{\small #1}}

\title{Two machine-checked refutations in the \Calc{} metatheory\\
  \large Lemma E.17 (Look completeness) and
  Theorem 4.6 (Completeness and optimality of subtyping)
  are false as stated}
\author{}
\date{}

\begin{document}
\maketitle

\begin{abstract}
  \noindent
  One subtyping configuration refutes both statements.  The declarative
  side of the configuration --- the solution, including the tag obligation
  it must discharge --- is a Coq term, not a prose argument; so is the
  non-existence of an algorithmic run.  No source program and no typing
  derivation is involved: the two statements are about subtyping and the
  refutation stays there.

  The defect is localised to \emph{one} rule of constraint solving,
  \soleqlab{FlexR}, which discards information the algorithm has already
  computed.  Handing the algorithm the instantiation that information
  determines makes the run exist --- also machine-checked.  A repair is
  suggested at the end, together with what still has to be checked before
  it is applied.

  All statements below are \texttt{Qed} and axiom-free
  (\texttt{Print Assumptions}: \emph{closed under the global context}).
\end{abstract}

% =====================================================================
\section{Summary}
% =====================================================================

\begin{center}\small
\begin{tabular}{@{}p{4.4cm}p{4.3cm}p{5.6cm}@{}}
  \toprule
  Statement & Where & Coq witness (\texttt{Qed}, axiom-free) \\
  \midrule
  Look completeness --- \textbf{false as stated}
    & \texttt{subtyping-proofs.tex:2730}
    & \coq{p6l\_look\_completeness\_fails} \\[.6ex]
  Completeness and optimality of subtyping --- \textbf{false as stated}
    & \texttt{algorithmic.tex:1447}
    & \coq{thm\_complete\_subtyping\_false} \\
  \midrule
  \emph{and the defect, localised}: the run \emph{does} exist at the
  correct instantiation
    & (a positive fact about the paper's own algorithm)
    & \coq{run\_exists\_at\_correct\_instantiation}\\
  \bottomrule
\end{tabular}
\end{center}

\noindent
Files:
\coq{frost/algo/cex/probe\_double\_occ\_run.v},
\coq{frost/algo/cex/cex\_thm\_complete\_subtyping.v},
\coq{frost/algo/cex/probe\_solve\_flexr.v}.

% =====================================================================
\section{The configuration}\label{sec:config}
% =====================================================================

Throughout, $\Theta_0$ holds exactly one \emph{undefined} flexible
variable, and
\[
  \Theta_0 = \hat{f}
  \qquad
  \N = \TY
  \qquad
  \M = \FN\TC
  \qquad
  P_0 = (\forall \gamma.\gamma\to\gamma) \to \hat{f}
  \qquad
  P   = \forall \beta.\beta\to\beta .
\]
The metasubstitution, the \sublab{ForallL} instantiation and the
counterpart of the declarative solution are
\[
  \theta = [\hat f \mapsto \tunit\to\tunit]
  \qquad
  \psi_a = \foralltag{\epsilon}\beta.\beta\to\beta
  \qquad
  \Psi_a = \{\epsilon\}
  \qquad
  Q = \inh{(\forall\beta.\beta\to\beta)} \to (\tunit\to\tunit).
\]

Two features of $P_0$ are what the refutation turns on, and it is worth
naming them now.
\begin{itemize}
\item Its \emph{domain} is a bare $\forall$.  An arrow-left is a
  \emph{consistency} position, and $\inh{\theta P_0}$ paints exactly the
  $\forall$s of $P_0$ inherited --- so this is the only position at which
  a guess's tag can be paid.
\item Its \emph{codomain} is a raw $\Theta_0$-flexible.  The algorithm has
  no inherited information there, while the declarative side reads
  $\theta\hat f = \tunit\to\tunit$.
\end{itemize}

% =====================================================================
\section{Completeness of subtyping is false}\label{sec:thm}
% =====================================================================

We recall the two statements as the paper gives them.

\begin{defpaper}[\normalfont\texttt{algorithmic.tex:1423}]
  A subtyping problem is a tuple $(\Theta_0;\N;\M;P_0;P)$ where $\Theta_0$
  is well-formed, $P_0$ and $P$ are well-formed in $\Theta_0$, and
  $\skmod[\N]{P_0}{\M}$.
  %
  A solution to it is a tuple $(\inc[\theta_1]{\Theta_0}{\Theta_1};P_1)$
  such that $\typsub[\N][\M][]{\Theta_1}{\syn{\theta_1 P}}{\bid{P_1}}$
  where $\Theta_1\vdash\inh{\theta_1 P_0}\fillsyn\bid{P_1}$.
\end{defpaper}

\begin{thmpaper}[\normalfont Completeness and optimality of subtyping,
    \texttt{algorithmic.tex:1447}]
  Given a subtyping problem $(\Theta_0;\N;\M;P_0;P)$, if there exists a
  solution, $\atypsub[\N][\M][P_0]{\Theta_0}{P}{P_1}{\Theta_1}$ gives its
  optimal solution $(\inc{\Theta_0}{\Theta_1};P_1)$.
\end{thmpaper}

\begin{claimc}\label{cl:thm}
  Theorem 4.6 is false at the configuration of \Cref{sec:config}.
\end{claimc}

\paragraph{(1) It is a subtyping problem.}
$\Theta_0$ is well-formed; $\Theta_0\vdash P_0$ and $\Theta_0\vdash P$;
and $\skmod[\TY]{P_0}{\FN\TC}$ holds, because $\skmod[\TY]{\cdot}{\TC}$
asks for a \emph{type} and the spine of $P_0$ ends in the type $\hat f$.
(Coq: \coq{$\Theta$0d\_wf}, \coq{P0d\_wf}, \coq{Pd\_wf}, \coq{P0d\_mode}.)

\paragraph{(2) It has a solution.}
Take $\theta_1=\theta$, $\Theta_1=\Theta_0$ and $P_1=Q$.  The fill
$\Theta_0\vdash\inh{\theta P_0}\fillsyn\bid{Q}$ holds
(\coq{fill\_d}), and the declarative derivation is
\[\small
\inferrule*[Lab=\sublab{ForallL}]
{
  \Theta_0\vdash^{\TY}\tags{\Psi_a}{\syn{\psi_a}}
  \\\\
  \inferrule*[Lab=\sublab{Arrow}]
  {
    \inferrule*[Lab=\sublab{ForallInh}]
    { \{\epsilon\}\mid\Theta_0 \vdash \beta\to\beta \consis
      \inh{\beta\to\beta} }
    { \{\epsilon\}\mid\Theta_0 \vdash \syn{\psi_a} \consis
      \inh{(\forall\beta.\beta\to\beta)} }
    \\
    \inferrule*[Lab=\sublab{ForallL}]
    { \typsub[\TY][\TC][\cdot]{\Theta_0}{\syn{\tunit\to\tunit}}
        {\bid{\tunit\to\tunit}} }
    { \typsub[\TY][\TC][\cdot]{\Theta_0}{\syn{\psi_a}}
        {\bid{\tunit\to\tunit}} }
  }
  { \typsub[\TY][\FN\TC][\Psi_a]{\Theta_0}
      {\syn{\psi_a}\to\syn{\psi_a}}
      {\bid{\inh{(\forall\beta.\beta\to\beta)}\to(\tunit\to\tunit)}} }
}
{ \typsub[\TY][\FN\TC][\cdot]{\Theta_0}{\syn{\theta P}}{\bid{Q}} }
\]
The tag index splits at \sublab{Arrow} as $\Psi_a = \{\epsilon\}\uplus\cdot$.
\emph{The guess's tag $\epsilon$ is paid} on the left, by the tagged
\sublab{ForallInh} against an \emph{inherited} $\forall$ --- the ``never
guess polymorphism out of thin air'' obligation, discharged rather than
assumed.  On the right the same guess is instantiated with the
\emph{monotype} $\tunit$, so no tag is owed there.
(Coq: \coq{solution\_d}, built from \coq{consis\_d}, \coq{refine\_dom\_d},
\coq{join\_cod\_d}; the statement quotes the paper's judgement
$\typsub[\N][\M][\Psi]{\Theta}{\cdot}{\cdot}$, not any auxiliary
mechanised variant.)

\paragraph{(3) Yet no algorithmic run exists.}
\[
  \nexists\, P_1\ \Theta_1.\quad
  \atypsub[\TY][\FN\TC][P_0]{\Theta_0}{P}{P_1}{\Theta_1}
\]
(Coq: \coq{probe\_no\_run}), so Theorem 4.6 fails --- a fortiori for
optimality.  The trace is forced; \Cref{sec:why} gives it.
\qed

% =====================================================================
\section{Look completeness is false}\label{sec:look}
% =====================================================================

Theorem 4.6's \sublab{ForallL} case (\texttt{subtyping-proofs.tex:3371})
proves itself by invoking the following lemma, and the same configuration
refutes it directly --- one lemma lower, and with no appeal to the
theorem.

\begin{lempaper}[\normalfont Look completeness,
    \texttt{subtyping-proofs.tex:2730}]
  Suppose
  $\Theta_0\vdash P_0$, $\skmod[\N]{P_0}{\M}$, $\erase{S}=P$ and
  $\inc[\theta]{\Theta_0}{\Theta}$.  If $\psi_a$ is a valid instantiation
  for declarative \sublab{ForallL} with
  \[
    \Theta\vdash^\N\tags{\Psi_a}{\syn{\psi_a}}
    \quad
    \Psi \text{ disjoint from } \Psi_a
    \quad
    \typsub[\N][\M][\Psi,\Psi_a]{\Theta}{
      (\syn{\theta S})[\syn{\psi_a}/\alpha]}{\bid{Q}}
    \quad
    \Theta\vdash\inh{\theta P_0}\fillsyn\bid{Q},
  \]
  then there exist $P_a$, $\Xi$ and $\theta_a$ such that
  $\looksub[\N][\M]{\alpha}{P}{P_0}{P_a}[\Xi]$ and
  $\Theta_0,\Xi\vdash P_a$, with
  $\inc[\theta_a]{\Theta_0,\Xi}{\Theta}$,
  $\inc[\theta\equiv\theta_a]{\Theta_0}{\Theta}$ and
  \[
    \Theta\vdash\precise{\theta_a P_a}{\erase{\psi_a}} .
  \]
\end{lempaper}

\begin{claimc}
  Lemma E.17 is false at the configuration of \Cref{sec:config},
  with $S = \foralltag{t}\beta.\beta\to\beta$ and $\Psi=\cdot$.
\end{claimc}

\paragraph{Every premise holds.}
$\Theta_0\vdash P_0$ (\coq{P0d\_wf}); $\skmod[\TY]{P_0}{\FN\TC}$
(\coq{P0d\_mode}); $\erase{S}=P$ by computation;
$\inc[\theta]{\Theta_0}{\Theta_0}$ (\coq{inc\_d}); the tag sorting
$\Theta_0\vdash^{\TY}\tags{\Psi_a}{\syn{\psi_a}}$ (\coq{$\psi$d\_tags}) with
$\Psi=\cdot$ trivially disjoint from $\Psi_a$; the declarative premise at
$\Psi,\Psi_a$ (\coq{look\_premise\_d}, again in the paper's judgement);
and the fill (\coq{fill\_d}).

\paragraph{The conclusion fails, for \emph{every} output.}
The look is deterministic here (\coq{probe\_extract\_det}), uniformly in
the marker and in every generated name --- the binder it produces is
vacuous, so it swallows them --- and its only output is
\[
  P_a \;=\; \forall\_.\,\hat f .
\]
The lemma then demands $\Theta_0\vdash\precise{\theta P_a}{\erase{\psi_a}}$,
i.e.
\[
  \precise{\forall\_.\,(\tunit\to\tunit)}{\forall\beta.\beta\to\beta} ,
\]
which is refuted: one arrow step leaves $\tunit$ against the rigid
$\beta$, and the only rule that could refine a unit to a variable requires
a \emph{defined} flexible, whereas $\Theta_0$ carries an undefined one.
(Coq: \coq{p6l\_look\_completeness\_fails}.)
\qed

% =====================================================================
\section{Why the algorithm fails: the marker occurs twice}\label{sec:why}
% =====================================================================

Algorithmic \silab{ForallL} picks a fresh $\alpha$ --- the \emph{marker}
--- opens the subject's binder with it, and calls the extraction judgement
\emph{indexed by} $\alpha$, which emits one constraint per
\emph{occurrence} of the marker.  Here the subject is
$\forall\beta.\beta\to\beta$, so the opened body is $\alpha\to\alpha$ and
the marker occurs \emph{twice}.  The two occurrences sit at different
kinds of position, so they take different branches of the collection:
\[
  \looksub[\FN\TC]{\alpha}{\alpha\to\alpha}{P_0}
    {\eqcons{\forall\gamma.\gamma\to\gamma},\;
     \subcons{\TC}{\hat f}}[\cdot]
\]
\begin{center}
\begin{tabular}{@{}llll@{}}
  \toprule
  occurrence & faces in $P_0$ & position & rule and constraint\\
  \midrule
  domain   & $\forall\gamma.\gamma\to\gamma$ & arrow-left: consistency
           & \lelab{Var}, $\eqcons{\forall\gamma.\gamma\to\gamma}$\\
  codomain & $\hat f$ & arrow-right: subtyping
           & \lslab{Var}, $\subcons{\TC}{\hat f}$\\
  \bottomrule
\end{tabular}
\end{center}
\noindent
(\coq{probe\_look\_det}: that list, deterministically.)  One marker, two
constraints --- so the solve has to \emph{combine} them, and that is where
the information is lost.  With a single occurrence there would be one
constraint and nothing to lose.

\paragraph{The forced trace.}
The solve threads the accumulated pattern through the constraint list,
starting from $\Hole$:
\[\bl
  \solveeq{\Hole}{\forall\gamma.\gamma\to\gamma}
          {\forall\gamma.\gamma\to\gamma}
  \hfill \soleqlab{GhostL} \\[.4ex]
  \solvesub{\forall\gamma.\gamma\to\gamma}{\TC}{\hat f}{\forall\_.\hat f}
  \hfill \solsublab{ForallL},\ \solsublab{Check},\ \soleqlab{FlexR} \\[.4ex]
  \guess[\TY]{\forall\_.\hat f} = \forall\_.\hat f
\el\]
The middle line is the defect.  \solsublab{ForallL} opens $\beta$ and
\solsublab{Check} hands $\beta\to\beta \consis \hat f$ to
\soleqlab{FlexR}, whose output \emph{is} the flexible $\hat f$; the close then
re-binds a \emph{vacuous} binder.  Continuing with that instantiation,
\silab{Arrow}'s domain consistency
$\forall\_.\hat f \consis \forall\gamma.\gamma\to\gamma$ pushes a fresh
rigid $\beta_1$, leaving the bare $\hat f$ against $\beta_1\to\beta_1$,
where the assignment walk meets $\beta_1$ with $\beta_1$ occurring --- and
is stuck (\coq{probe\_solve\_det}, \coq{probe\_continuation\_stuck}).

\paragraph{One rule treats ``nothing is known here'' two different ways.}
Put the three relevant rows side by side:
\[
\inferrule*[Lab=\soleqlab{FlexR}]
{\nonflex{P}}
{\solveeq{P}{\hat\alpha}{\hat\alpha}}
\qquad
\inferrule*[Lab=\soleqlab{VacuousR}]
{~}
{\solveeq{P}{\Hole}{P}}
\qquad
\inferrule*[Lab=\solsublab{Vacuous}]
{~}
{\solvesub{P}{\M}{\Hole}{P}}
\]
A raw flexible and a ghost are the same epistemic situation: nothing is
known at that position.  The two ghost rows keep the pattern already
accumulated; \soleqlab{FlexR} replaces it by the flexible.  The collection
itself takes the other side --- on the equality branch, \lelab{FlexR}
facing a flexible input invents a \emph{fresh} flexible and emits an
equality against it, rather than letting the input's flexible win.

\paragraph{And the loss is the whole defect.}
Hand \silab{ForallL} the instantiation the discarded equality determines,
$\forall\beta.\beta\to\beta$ (that is $\erase{\psi_a}$), and the run
exists:
\[
  \exists\, Q\ \Theta_1.\quad
  \atypsub[\TY][\FN\TC][P_0]{\Theta_0}
    {(\forall\beta.\beta\to\beta)\to(\forall\beta.\beta\to\beta)}
    {Q}{\Theta_1}
\]
(Coq: \coq{run\_exists\_at\_correct\_instantiation}).  The run it builds is
the expected one: the domain by \sublab{ForallInh} at a fresh rigid, then
\sublab{Arrow} of two rigids; the codomain
$\forall\beta.\beta\to\beta \subtype_{\TC} \hat f$ by \silab{ForallL},
whose collection takes \lslab{Check} and then \lelab{FlexR} --- inventing a
fresh $\hat\alpha_1$ --- and whose continuation
$\hat\alpha_1\to\hat\alpha_1 \subtype_{\TC} \hat f$ goes through
\silab{Check}, \eqlab{FlexR} and the assignment, setting
$\hat f \coloneqq \hat\alpha_1\to\hat\alpha_1$.  That solution is
\emph{more general} than the declarative one: $\theta\hat f =
\tunit\to\tunit$ factors through it by $\hat\alpha_1\mapsto\tunit$.

So every rule the run needs already exists and already works.  What the
algorithm cannot do is \emph{choose} $\forall\beta.\beta\to\beta$.

% =====================================================================
\section{Suggested repair}
% =====================================================================

\soleqlab{FlexR} should yield the accumulated pattern, as its ghost twin
does; the flexible is recorded in $\Theta$ by the assignment that follows,
not by the solve's output:
\[
\inferrule*[Lab=\soleqlab{FlexR}$'$]
{\nonflex{P}}
{\solveeq{P}{\hat\alpha}{P}}
\]
The case that makes the present reading useful --- $P = \Hole$, where
returning $\hat\alpha$ lets the guess \emph{share} the flexible instead of
inventing a fresh one --- is already covered by \soleqlab{GhostL}
($\solveeq{\Hole}{Q}{Q}$), so the two readings differ only where the
pattern already carries information, which is exactly where the present
one destroys it.

\medskip\noindent
\textbf{What has not been checked}, and should be before the change is
made: soundness of subtyping under the changed row, and the two collection
/ solving completeness lemmas that the subtyping completeness proof rests
on.  The mechanisation therefore does \emph{not} apply the change; it
reports it.

% =====================================================================
\section*{Index of Coq names}
% =====================================================================

\noindent\small
\begin{tabular}{@{}p{6.2cm}p{9.2cm}@{}}
  \toprule
  \coq{thm\_complete\_subtyping\_false}
    & Theorem 4.6 refuted (\Cref{cl:thm})\\
  \coq{solution\_d}, \coq{look\_premise\_d}
    & the declarative solution / premise, in the paper's judgement\\
  \coq{consis\_d}, \coq{join\_cod\_d}, \coq{refine\_dom\_d},
    \coq{fill\_d} & its leaves; \coq{consis\_d} is where $\epsilon$ is paid\\
  \coq{p6l\_look\_completeness\_fails}
    & Lemma E.17 refuted\\
  \coq{probe\_look\_det}, \coq{probe\_solve\_det},
    \coq{probe\_extract\_det} & the collection / solve / guess are forced\\
  \coq{probe\_continuation\_stuck}, \coq{probe\_no\_run}
    & no algorithmic run exists\\
  \coq{run\_exists\_at\_correct\_instantiation}
    & \dots but one does at $\erase{\psi_a}$\\
  \bottomrule
\end{tabular}

\end{document}
